\subsubsection{A regularization defect along a complex disk}
\label{sec:affine-contact-defect}
\index{Contact set!regularization defect}

Raw contact with a constant obstacle can hold for every positive growth parameter while regularization raises the value of the envelope at a point of the compact set. In the following example, a complex disk through the point, transverse to the planar direction, prevents any local plurisubharmonic function from separating the point from the rest of the compact set, but it does not remove this defect.

In the complex plane, let $D$ be the open unit disk and put
\[
 p_j=2+2^{-j},\qquad r_j=\exp(-16\,2^j),\qquad
 E=\overline D\cup\{2\}\cup\bigcup_{j\geq1}\overline D(p_j,r_j).
\]

\Needspace{17\baselineskip}
\begin{lemma}\label{lem:planar-contact-barrier}
The set $E$ is compact, and every neighborhood of each of its points contains an open disk contained in $E$. The series
\[
 h(\zeta)=\sum_{j\geq1}2^{-j}\log|\zeta-p_j|+2\log2
\]
defines a subharmonic function on $\mathbb C$ of logarithmic growth at most one; it satisfies $h(2)=0$ and is at most $-16+2\log2$ on every small disk $\overline D(p_j,r_j)$. The function
\[
 v(\zeta)=\frac{8\log^+|\zeta|+h(\zeta)-4\log2}{9}
\]
is subharmonic of logarithmic growth at most one, is nonpositive on $E\setminus\{2\}$, and satisfies $v(2)=4\log2/9$.
\end{lemma}
\begin{proof}
Since $r_j<2^{-j-3}$, the small disks are pairwise disjoint and disjoint from $\overline D$, and they accumulate only at $2$. This proves compactness. Every neighborhood of $2$ contains one of the small disks, and every other point of $E$ lies in the closure of the interior of one of the disks.

Fix $R>0$. On the disk $\{|\zeta|<R\}$ every term satisfies $\log|\zeta-p_j|<\log(R+3)$, so the partial sums of $\sum_j2^{-j}(\log|\zeta-p_j|-\log(R+3))$ decrease there. Their limit is therefore subharmonic or identically $-\infty$ on this disk, and it is finite at $\zeta=0$. Since $\sum_j2^{-j}=1$, this limit is $h-2\log2-\log(R+3)$, so $h$ is subharmonic on $\mathbb C$. The bound $|\zeta-p_j|\leq|\zeta|+3$ gives $h(\zeta)\leq\log(|\zeta|+3)+2\log2$, hence the growth bound, and $\sum_jj2^{-j}=2$ gives $h(2)=0$.

On the $j$th small disk, all distances to the other centers are less than one, so the corresponding terms are negative, and
\[
 h\leq2^{-j}\log r_j+2\log2=-16+2\log2.
\]
On $\overline D$ all distances to the centers are less than four, hence $h\leq4\log2$ and $v\leq0$. On each small disk $|\zeta|<4$, so $9v\leq14\log2-16<0$. The growth bound for $v$ and its value at $2$ follow directly from the formula.
\end{proof}

Now let $n\geq2$ and $K=E\times\overline D^{\,n-1}\subset\mathbb C^n$. For a non-pluripolar compact set $L\subset\mathbb C^n$, write $V_{\delta,L}$ and $F_{\delta,L}$ for the raw envelope with growth allowance $\delta$ and zero obstacle on $L$, and for its regularization; thus $F_{\delta,L}=F_{\delta,L,0}$ in the notation of \cref{sec:affine-local-regularity}.

\Needspace{19\baselineskip}
\begin{proposition}\label{prop:affine-transverse-contact-defect}
For every point of $K$ and every neighborhood $W$ of it, the set $K\cap W$ is non-pluripolar. For $\delta>0$ and $a=(2,w)$ with $w\in\overline D^{\,n-1}$,
\begin{equation}\label{eq:affine-contact-defect}
 V_{\delta,K}|_K=0,\qquad
 F_{\delta,K}(a)\geq\frac{4\delta}{9}\log2>0.
\end{equation}
For every $r>0$, the same lower bound holds for $F_{\delta,K\cap\overline B(a,r)}(a)$. At $a_0=(2,0,\ldots,0)$, every plurisubharmonic function $f$ defined near $a_0$ satisfies
\begin{equation}\label{eq:transverse-psh-no-jump}
 f(a_0)\leq\limsup_{K\setminus\{a_0\}\ni z\to a_0}f(z).
\end{equation}
Thus no plurisubharmonic function near $a_0$ separates $a_0$ from $K\setminus\{a_0\}$, that is, $K$ is not plurithin at $a_0$. Nevertheless, the regularized envelope does not touch the obstacle at $a_0$ for any positive growth parameter.
\end{proposition}
\index{Envelope!non-thin contact defect}
\begin{proof}
By the preceding lemma, every neighborhood of a point of $K$ contains a nonempty open subset of $K$, namely a product of open disks. This proves the first assertion, which applies in particular to each $K\cap\overline B(a,r)$, $r>0$.

The plurisubharmonic function $z\mapsto v(z_1)$ belongs to $\mathcal L_1$ and is nonpositive on $K$ outside the pluripolar set $K\cap\{z_1=2\}$. By \cref{cor:affine-polar-obstacle-invariance} with zero obstacle, $v(z_1)\leq F_{1,K}(z)$ for all $z$. The value $v(2)=4\log2/9$ gives the lower bound for $\delta=1$. Since $0$ is a competitor, $V_{\delta,K}|_K=0$, and scaling of competitors gives $F_{\delta,K}=\delta F_{1,K}$. Replacing $K$ by a smaller compact set only increases the envelope, which proves the local assertion.

Finally, let $f$ be plurisubharmonic near $a_0$, and restrict it to the complex line through $a_0$ in which only the second coordinate varies. Small circles in this line centered at $a_0$ lie in $K\setminus\{a_0\}$, and the submean inequality on these circles gives \eqref{eq:transverse-psh-no-jump}.
\end{proof}
